%% 2i2d-treillis-ferme — ferme en treillis à membrures parallèles : chaque barre
%% ne travaille qu'en traction ou en compression. Le triangle est indéformable,
%% le quadrilatère articulé ne l'est pas.
\documentclass[tikz,border=4pt]{standalone}
\usepackage{liaisons}
\begin{document}
\begin{tikzpicture}[x=1cm,y=1cm,
  barre/.style={line width=0.9pt, solideB, line cap=round},
  tendue/.style={line width=3.4pt, solideA!45, line cap=round},
  comprimee/.style={line width=3.4pt, solideB!30, line cap=round},
  force/.style={-{Stealth[length=5pt]}, line width=1pt, solideA},
  etiq/.style={font=\small, inner sep=2pt},
  petit/.style={font=\scriptsize, inner sep=1.5pt}]

\newcommand\batihoriz[3]{%
  \draw[line width=1pt, solideE] (#1,#3) -- (#2,#3);
  \begin{scope}
    \clip (#1,#3-0.24) rectangle (#2,#3);
    \foreach \hx in {-16,-15.76,...,16} { \draw[line width=0.5pt, solideE!75] (\hx,#3) -- (\hx-0.28,#3-0.28); }
  \end{scope}}
\newcommand\appuitri[3]{%
  \draw[line width=1pt, draw=solideE, fill=solideE!12, line join=round]
    (#1,#2) -- (#1-0.28,#3) -- (#1+0.28,#3) -- cycle;}

%% --- surlignage des efforts (tracé sous les barres) --------------------------
\draw[tendue]    (0.6,1.0) -- (7.4,1.0);                     % membrure inférieure
\draw[tendue]    (0.6,1.0) -- (2.3,2.4)  (2.3,1.0) -- (4.0,2.4)
                 (5.7,2.4) -- (4.0,1.0)  (7.4,2.4) -- (5.7,1.0);   % diagonales
\draw[comprimee] (0.6,2.4) -- (7.4,2.4);                     % membrure supérieure
\foreach \x in {0.6,2.3,4.0,5.7,7.4} { \draw[comprimee] (\x,1.0) -- (\x,2.4); }

%% --- les barres ----------------------------------------------------------------
\draw[barre] (0.6,2.4) -- (7.4,2.4);
\draw[barre] (0.6,1.0) -- (7.4,1.0);
\foreach \x in {0.6,2.3,4.0,5.7,7.4} { \draw[barre] (\x,1.0) -- (\x,2.4); }
\draw[barre] (0.6,1.0) -- (2.3,2.4)  (2.3,1.0) -- (4.0,2.4)
             (5.7,2.4) -- (4.0,1.0)  (7.4,2.4) -- (5.7,1.0);

%% --- les nœuds -------------------------------------------------------------------
\foreach \x in {0.6,2.3,4.0,5.7,7.4} { \fill (\x,1.0) circle (1.3pt); \fill (\x,2.4) circle (1.3pt); }

%% --- les charges sur les nœuds supérieurs -------------------------------------------
\foreach \x in {2.3,4.0,5.7} { \draw[force] (\x,3.3) -- (\x,2.48); }

%% --- les appuis -----------------------------------------------------------------------
\appuitri{0.6}{0.97}{0.45}
\batihoriz{0.15}{1.05}{0.45}
\appuitri{7.4}{0.97}{0.52}
\draw[line width=0.8pt, draw=solideE, fill=white] (7.24,0.41) circle (0.1);
\draw[line width=0.8pt, draw=solideE, fill=white] (7.56,0.41) circle (0.1);
\batihoriz{6.95}{7.85}{0.31}

%% --- légende du code couleur ---------------------------------------------------------
\draw[tendue]    (4.40,4.42) -- (4.90,4.42);
\node[petit, anchor=west, text=solideA] at (4.98,4.42) {barre \textbf{tendue}};
\draw[comprimee] (4.40,3.88) -- (4.90,3.88);
\node[petit, anchor=west, text=solideB] at (4.98,3.88) {barre \textbf{comprimée}};

%% --- encart : triangle indéformable / quadrilatère déformable -------------------------
\draw[rounded corners=2.5pt, draw=black!35, fill=black!3] (0.1,3.42) rectangle (3.5,4.78);
\draw[line width=0.9pt, solideE, line join=round] (0.45,3.87) -- (1.15,3.87) -- (0.80,4.47) -- cycle;
\foreach \p in {(0.45,3.87),(1.15,3.87),(0.80,4.47)} { \fill \p circle (1.2pt); }
\node[petit, anchor=north, text=solideE] at (0.80,3.84) {indéformable};
\draw[line width=0.9pt, solideA, line join=round]
  (2.15,3.87) -- (2.85,3.87) -- (3.05,4.47) -- (2.35,4.47) -- cycle;
\foreach \p in {(2.15,3.87),(2.85,3.87),(3.05,4.47),(2.35,4.47)} { \fill \p circle (1.2pt); }
\node[petit, anchor=north, text=solideA] at (2.60,3.84) {déformable};
\end{tikzpicture}
\end{document}
